% ===========================================================================
% 03 - Technische Anforderungen (TA) im Code
% ===========================================================================

% --- TA-01 -----------------------------------------------------------------
\begin{frame}[fragile]{TA-01 -- Spring Boot}
  \tabadge{TA-01}
  \begin{columns}[T,onlytextwidth]
    \begin{column}{0.52\textwidth}
      \begin{lstlisting}[style=java]
@EnableAsync
@SpringBootApplication
public class FluidAnalysisServiceApplication {
    public static void main(String[] args) {
        SpringApplication.run(
            FluidAnalysisServiceApplication.class, args);
    }
}
      \end{lstlisting}
      \srcfile{services/fluid-analysis-service/.../FluidAnalysisServiceApplication.java}
    \end{column}
    \begin{column}{0.44\textwidth}
      \textbf{Erläuterung}
      \begin{itemize}
        \item \code{@SpringBootApplication} aktiviert Auto-Configuration, Component Scan und die eingebettete Servlet-Engine.
        \item \code{@EnableAsync} ist die Grundlage für die asynchrone Ausführung (QR-08).
        \item Identisches Muster in allen sechs Services.
      \end{itemize}
      \vspace{0.2cm}
      \begin{block}{Abhängigkeit}
        \small \code{pom.xml}: \code{spring-boot-starter-web}
      \end{block}
    \end{column}
  \end{columns}
\end{frame}

% --- TA-02 -----------------------------------------------------------------
\begin{frame}[fragile]{TA-02 -- Docker (Container pro Service)}
  \tabadge{TA-02}
  \begin{columns}[T,onlytextwidth]
    \begin{column}{0.54\textwidth}
      \begin{lstlisting}[style=docker]
FROM maven:3.9.16-eclipse-temurin-21 AS build
WORKDIR /workspace
COPY pom.xml ./
COPY services ./services
RUN mvn -q -pl services/fluid-analysis-service \
    -am clean package -DskipTests

FROM eclipse-temurin:21-jre
WORKDIR /app
COPY --from=build /workspace/services/fluid-analysis-service/target/fluid-analysis-service-0.1.0-SNAPSHOT.jar app.jar
EXPOSE 8083
ENTRYPOINT ["java", "-jar", "/app/app.jar"]
      \end{lstlisting}
      \srcfile{services/fluid-analysis-service/Dockerfile}
    \end{column}
    \begin{column}{0.42\textwidth}
      \textbf{Erläuterung}
      \begin{itemize}
        \item \textbf{Multi-Stage-Build}: Build-Umgebung (Maven + JDK) getrennt vom schlanken Runtime-Image (nur JRE).
        \item Jeder Service erhält ein eigenes Image und einen eigenen Port.
        \item Grundlage für Independent Deployability (QR-03, ADR-05).
      \end{itemize}
      \vspace{0.2cm}
      \begin{block}{Nachweis}
        \small 6 $\times$ \code{Dockerfile} -- eines pro Service.
      \end{block}
    \end{column}
  \end{columns}
\end{frame}

% --- TA-03 -----------------------------------------------------------------
\begin{frame}[fragile]{TA-03 -- Docker Compose}
  \tabadge{TA-03}
  \begin{columns}[T,onlytextwidth]
    \begin{column}{0.54\textwidth}
      \begin{lstlisting}[style=yaml]
services:
  configuration-service:
    image: ysirin2s/seka-wirschaffendas:configuration-service-v${VERSION}
    ports:
      - "8081:8081"
    volumes:
      - configuration-data:/app/data
  analysis-management-service:
    image: ysirin2s/seka-wirschaffendas:analysis-management-service-v${VERSION}
    ports:
      - "8082:8082"
    environment:
      CONFIGURATION_SERVICE_URL: http://configuration-service:8081
      FLUID_ANALYSIS_SERVICE_URL: http://fluid-analysis-service:8083
    volumes:
      - analysis-data:/app/data
      \end{lstlisting}
      \srcfile{docker-compose.yml}
    \end{column}
    \begin{column}{0.42\textwidth}
      \textbf{Erläuterung}
      \begin{itemize}
        \item Compose definiert alle sechs Services, Ports, Umgebungsvariablen und getrennte Volumes.
        \item Services finden sich über \textbf{Service-Namen} im Netzwerk \code{wirschaffendas-network} -- keine fest verdrahteten IPs.
        \item \code{depends_on} steuert die Startreihenfolge.
      \end{itemize}
      \begin{block}{Start}
        \small \code{docker compose up --build}
      \end{block}
    \end{column}
  \end{columns}
\end{frame}

% --- TA-04 -----------------------------------------------------------------
\begin{frame}[fragile]{TA-04 -- Resilience4j: Circuit Breaker}
  \tabadge{TA-04}
  \begin{columns}[T,onlytextwidth]
    \begin{column}{0.56\textwidth}
      \begin{lstlisting}[style=java]
@CircuitBreaker(name = "nextService", fallbackMethod = "fallback")
public void startNext(AnalysisCommand command, String currentResult) {
    var results = new ArrayList<>(command.previousResults());
    results.add(Map.of("algorithm", "FLUID", "result", currentResult));
    AnalysisCommand nextCommand = new AnalysisCommand(
        command.analysisId(), command.configuration(), results);
    client.post().uri("/internal/analyses")
        .body(nextCommand).retrieve().toBodilessEntity();
}
private void fallback(AnalysisCommand command,
                      String currentResult, Throwable throwable) {
    managementClient.reportStatus(command.analysisId(), "THERMAL",
        "FAILED", "thermal-analysis-service unavailable");
}
      \end{lstlisting}
      \srcfile{services/fluid-analysis-service/.../infrastructure/NextServiceClient.java}
    \end{column}
    \begin{column}{0.4\textwidth}
      \textbf{Erläuterung}
      \begin{itemize}
        \item Der Breaker fängt Verbindungsfehler ab.
        \item Statt Fehlerkaskade läuft der \textbf{Fallback}: markiert den Schritt als \code{FAILED}.
        \item Der aufrufende Service bleibt erreichbar.
        \item Nach Wartezeit $\rightarrow$ Half-Open $\rightarrow$ Retry möglich.
      \end{itemize}

      \vspace{-0.1cm}
      \begin{block}{Wirkung}
        \small Isolation of Failures -- der Fehler bleibt lokal (ADR-04).
      \end{block}
    \end{column}
  \end{columns}
\end{frame}

\begin{frame}[fragile]{TA-04 -- Circuit Breaker: Konfiguration}
  \tabadge{TA-04}
  \begin{columns}[T,onlytextwidth]
    \begin{column}{0.5\textwidth}
      \begin{lstlisting}[style=yaml]
resilience4j:
  circuitbreaker:
    instances:
      nextService:
        sliding-window-size: 2
        minimum-number-of-calls: 1
        permitted-number-of-calls-in-half-open-state: 1
        failure-rate-threshold: 50
        wait-duration-in-open-state: 10s
        automatic-transition-from-open-to-half-open-enabled: true
      \end{lstlisting}
      \srcfile{services/fluid-analysis-service/src/main/resources/application.yml}
    \end{column}
    \begin{column}{0.46\textwidth}
      \textbf{Zustandsmodell}
      \begin{center}
      \begin{tikzpicture}[
        node distance=9mm,
        every node/.style={font=\scriptsize},
        st/.style={draw=darkblue,rounded corners,fill=lightblue,minimum width=15mm,minimum height=6mm,align=center},
        op/.style={draw=warnorange,rounded corners,fill=warnorange!10,minimum width=15mm,minimum height=6mm,align=center}
      ]
      \node[st] (c) {Closed};
      \node[op,right=of c] (o) {Open};
      \node[st,right=of o] (h) {Half-Open};
      \draw[-{Latex},thick] (c) -- node[above]{Fehler} (o);
      \draw[-{Latex},thick] (o) to[bend left=20] node[above]{10s} (h);
      \draw[-{Latex},thick] (h) to[bend left=20] node[below]{Fehler} (o);
      \draw[-{Latex},thick] (h) to[bend left=45] node[below]{Erfolg} (c);
      \end{tikzpicture}
      \end{center}

      \begin{block}{Health-Endpoint}
        \small \code{/actuator/health} zeigt den Breaker-Zustand.
      \end{block}
    \end{column}
  \end{columns}
\end{frame}

% --- TA-05 -----------------------------------------------------------------
\begin{frame}{TA-05 -- Apache Kafka (bewusst nicht umgesetzt)}
  \tabadge{TA-05}
  \begin{columns}[T,onlytextwidth]
    \begin{column}{0.52\textwidth}
      \textbf{Anforderung}
      \begin{itemize}
        \item Priorität \textbf{COULD} -- nicht Teil des Pflichtumfangs.
        \item Kafka könnte Statusnachrichten asynchron übertragen.
      \end{itemize}

      \vspace{0.3cm}
      \textbf{Entscheidung}
      \begin{itemize}
        \item Kein Kafka im Code.
        \item Statusmeldungen laufen synchron per REST-Callback (FR-09).
      \end{itemize}
    \end{column}
    \begin{column}{0.44\textwidth}
      \begin{block}{Begründung (ADR-02)}
        \small Für den kleinen PoC ist REST einfach nachvollziehbar, testbar und mit der Aufgabenstellung kompatibel. Kafka bleibt bewusst außerhalb des Kernumfangs.
      \end{block}
      \vspace{0.2cm}
      \begin{block}{Trade-off}
        \small REST erzeugt temporale Kopplung. Für eine produktive Variante wäre asynchrones Messaging zu prüfen.
      \end{block}
    \end{column}
  \end{columns}
\end{frame}
